%%%%%%%%%%%%%%%%%%%%%%%%%%%%%%%%%%%%%%%%%%%%%%%%%%%%%%%%%%%%%%%%%%%%%%%%%%%%%%%
% Generative modeling: VAEs and diffusion models
%%%%%%%%%%%%%%%%%%%%%%%%%%%%%%%%%%%%%%%%%%%%%%%%%%%%%%%%%%%%%%%%%%%%%%%%%%%%%%%

\subsection{Variational Inference and the Variational Autoencoder}
\label{sec:vae}

A \emph{generative model} aims to produce samples $x\in\cX$ whose distribution matches a data distribution $p_{\text{data}}(x)$.
A common approach introduces a latent variable $Z$ with a simple prior and defines a flexible conditional model for $X$ given $Z$.

\paragraph{Latent-variable model.}
Choose a prior $p(z)$ (e.g., $\cN(0,I)$), and a conditional decoder family $p_\theta(x\mid z)$ parametrized by $\theta$.
This defines a joint model
\begin{equation}
    p_\theta(x,z)=p(z)\,p_\theta(x\mid z)
\end{equation}
and a marginal likelihood
\begin{equation}
    p_\theta(x)=\int p_\theta(x,z)\,dz.
\end{equation}

\paragraph{Maximum likelihood training.}
Given i.i.d. data $x_1,\dots,x_N\sim p_{\text{data}}$, maximum likelihood estimation (MLE) solves
\begin{equation}
    \theta^* \in \arg\max_{\theta}\; \frac{1}{N}\sum_{i=1}^N \log p_\theta(x_i)
    \quad\Longleftrightarrow\quad
    \theta^* \in \arg\min_{\theta}\; \frac{1}{N}\sum_{i=1}^N -\log p_\theta(x_i).
\end{equation}
The integral in $p_\theta(x)$ is typically intractable, which motivates variational inference.

\paragraph{Variational posterior.}
Introduce an approximate posterior (encoder) family $q_\phi(z\mid x)$.
Then
\begin{equation}
\label{eq:vae-elbo-jensen}
\begin{aligned}
\log p_\theta(x)
&= \log \int p_\theta(x,z)\,dz
= \log \int q_\phi(z\mid x)\,\frac{p_\theta(x,z)}{q_\phi(z\mid x)}\,dz\\
&= \log\, \Esub{z\sim q_\phi(\cdot\mid x)}{\frac{p_\theta(x,z)}{q_\phi(z\mid x)}}
\;\ge\; \Esub{z\sim q_\phi(\cdot\mid x)}{\log\frac{p_\theta(x,z)}{q_\phi(z\mid x)}}.
\end{aligned}
\end{equation}
The inequality is Jensen's inequality (concavity of $\log$).
The right-hand side is the \emph{evidence lower bound} (ELBO).

\begin{definition}[ELBO]
For a given data point $x$, define
\begin{equation}
\label{eq:elbo-def}
    \mathcal{L}_{\text{ELBO}}(\theta,\phi;x)
    \coloneqq \Esub{z\sim q_\phi(\cdot\mid x)}{\log p_\theta(x\mid z)}
    - \KL{q_\phi(z\mid x)}{p(z)}.
\end{equation}
\end{definition}

\paragraph{Interpretation.}
The ELBO has two terms:
\begin{itemize}[leftmargin=2em]
    \item \textbf{Reconstruction / likelihood term:} encourage $p_\theta(x\mid z)$ to assign high probability to the data when $z\sim q_\phi(z\mid x)$.

    \item \textbf{Regularization term:} push the approximate posterior $q_\phi(z\mid x)$ toward the prior $p(z)$.
    This prevents the encoder from using an arbitrarily complex latent code and is essential for sampling.
\end{itemize}

\paragraph{Why ELBO is a lower bound.}
A useful identity is
\begin{equation}
\label{eq:elbo-gap}
    \log p_\theta(x)
    = \mathcal{L}_{\text{ELBO}}(\theta,\phi;x)
    + \KL{q_\phi(z\mid x)}{p_\theta(z\mid x)}.
\end{equation}
Since KL divergence is nonnegative, the ELBO is always a lower bound on the log-likelihood.
Equality holds iff $q_\phi(z\mid x)=p_\theta(z\mid x)$.

\paragraph{Common Gaussian parameterization.}
VAEs often take
\(
    q_\phi(z\mid x)=\cN\bigl(\mu_\phi(x),\,\mathrm{diag}(\sigma_\phi(x)^2)\bigr)
\)
with a neural network outputting $\mu_\phi(x)$ and $\sigma_\phi(x)$.
For continuous data (e.g., images scaled to $[0,1]$), one may use a Gaussian or logistic likelihood for $p_\theta(x\mid z)$.

\begin{remark}[Bits vs. nats]
In these notes, $\log$ is base $2$ (bits). In much of the VAE literature, $\ln$ is used (nats).
The objectives differ by a constant factor of $\log_2(e)$.
\end{remark}

\subsection{Diffusion Models as Hierarchical VAEs}
\label{sec:diffusion}

A diffusion model defines a \emph{fixed} Markov chain that gradually corrupts data into Gaussian noise (forward process), and learns to reverse this corruption (reverse process).
It can be viewed as a latent-variable model with latents $(X_1,\dots,X_T)$ and observed variable $X_0$.

\subsubsection*{Forward (noising) process}

Fix a variance schedule $\{\beta_t\}_{t=1}^T$ with $0<\beta_t<1$ and define $\alpha_t\coloneqq 1-\beta_t$ and $\bar\alpha_t\coloneqq \prod_{s=1}^t \alpha_s$.
The forward process is
\begin{equation}
    q(x_{1:T}\mid x_0)=\prod_{t=1}^T q(x_t\mid x_{t-1}),
    \qquad
    q(x_t\mid x_{t-1}) = \cN\bigl(\sqrt{\alpha_t}\,x_{t-1},\,\beta_t I\bigr).
\end{equation}
One can show by iterating Gaussians that
\begin{equation}
\label{eq:xt-given-x0}
    q(x_t\mid x_0) = \cN\bigl(\sqrt{\bar\alpha_t}\,x_0,\,(1-\bar\alpha_t)I\bigr),
\end{equation}
so that $X_t$ is a noisy version of $X_0$ with signal-to-noise ratio controlled by $\bar\alpha_t$.
As $t\to T$ with $\bar\alpha_T\approx 0$, $X_T$ is nearly standard normal.

\subsubsection*{Reverse (denoising) model}

The learned generative process runs backward:
\begin{equation}
    p_\theta(x_{0:T}) = p(x_T)\prod_{t=1}^T p_\theta(x_{t-1}\mid x_t),
    \qquad p(x_T)=\cN(0,I).
\end{equation}
A common choice is a Gaussian reverse kernel
\begin{equation}
    p_\theta(x_{t-1}\mid x_t)=\cN\bigl(\mu_\theta(x_t,t),\,\Sigma_t\bigr),
\end{equation}
where $\Sigma_t$ is fixed (or learned) and a neural network parametrizes the mean.

\subsubsection*{ELBO decomposition}

As with VAEs, we can lower-bound the data log-likelihood $\log p_\theta(x_0)$ by an ELBO built on the forward chain $q$.
A standard calculation yields
\begin{equation}
\label{eq:diffusion-elbo}
\log p_\theta(x_0)
\ge
\Esub{q(x_{1:T}\mid x_0)}{\log p_\theta(x_{0:T}) - \log q(x_{1:T}\mid x_0)}
= L_0 - L_T - \sum_{t=1}^{T-1} L_t,
\end{equation}
where (up to constants)
\begin{align}
    L_0 &\equiv -\Esub{q(x_1\mid x_0)}{\log p_\theta(x_0\mid x_1)},\\
    L_T &\equiv \KL{q(x_T\mid x_0)}{p(x_T)},\\
    L_t &\equiv \Esub{q(x_t\mid x_0)}{\KL{q(x_{t-1}\mid x_t,x_0)}{p_\theta(x_{t-1}\mid x_t)}},\qquad t=1,\dots,T-1.
\end{align}
The dominant learning signal is matching the true reverse conditional $q(x_{t-1}\mid x_t,x_0)$.

\subsubsection*{Closed form of the true reverse conditional}

Because the forward chain is linear-Gaussian, the conditional $q(x_{t-1}\mid x_t,x_0)$ is Gaussian and can be computed in closed form.
One convenient form is
\begin{equation}
\label{eq:true-reverse}
q(x_{t-1}\mid x_t,x_0)
=\cN\bigl(\mu_q(x_t,x_0),\,\tilde\beta_t I\bigr),
\end{equation}
with
\begin{equation}
\mu_q(x_t,x_0)=
\frac{\sqrt{\bar\alpha_{t-1}}\,\beta_t}{1-\bar\alpha_t}\,x_0
+\frac{\sqrt{\alpha_t}(1-\bar\alpha_{t-1})}{1-\bar\alpha_t}\,x_t,
\qquad
\tilde\beta_t\coloneqq \frac{1-\bar\alpha_{t-1}}{1-\bar\alpha_t}\,\beta_t.
\end{equation}
Thus, learning the reverse kernel reduces to estimating $x_0$ (or equivalently the noise) given $x_t$.

\subsection{Equivalent Parameterizations: Predicting \texorpdfstring{$x_0$}{x0}, the Score, or the Noise}
\label{sec:diffusion-parameterizations}

Equation~\eqref{eq:xt-given-x0} implies a convenient reparameterization:
\begin{equation}
    X_t = \sqrt{\bar\alpha_t}\,X_0 + \sqrt{1-\bar\alpha_t}\,\varepsilon,
    \qquad \varepsilon\sim \cN(0,I).
\end{equation}
Given a sample $(x_0,t,\varepsilon)$, we can generate $x_t$ exactly.
Training is then supervised regression.

\subsubsection*{Predicting $x_0$}

A network can be trained to output an estimate $\widehat x_0 = x_{0,\theta}(x_t,t)$.
In the simplest (unweighted) form one minimizes
\begin{equation}
    \Esub{(x_0,t,\varepsilon)}{\|x_0 - x_{0,\theta}(x_t,t)\|^2},
    \qquad x_t=\sqrt{\bar\alpha_t}x_0+\sqrt{1-\bar\alpha_t}\,\varepsilon.
\end{equation}
The optimal predictor under squared error is the conditional mean $\E{X_0\mid X_t=x_t}$.

\subsubsection*{Predicting the noise $\varepsilon$}

Solving for $x_0$ gives
\(
    x_0 = \frac{x_t-\sqrt{1-\bar\alpha_t}\,\varepsilon}{\sqrt{\bar\alpha_t}}.
\)
Substituting into $\mu_q(x_t,x_0)$ shows that the reverse mean can be written as an affine function of $x_t$ and $\varepsilon$.
Therefore, instead of predicting $x_0$, one can predict the noise $\varepsilon$ via a network $\varepsilon_\theta(x_t,t)$ and train using
\begin{equation}
\label{eq:eps-loss}
    \Esub{(x_0,t,\varepsilon)}{\|\varepsilon - \varepsilon_\theta(x_t,t)\|^2}.
\end{equation}
This is the most common parameterization in practice.

\subsubsection*{Predicting the score}

The \emph{score function} of the marginal $p_t(x)\coloneqq q(x_t=x)$ is
\(\nabla_x\log p_t(x)\).
Tweedie's formula relates the posterior mean to the score:
\begin{equation}
    \E{X_0\mid X_t=x_t}
    = \frac{1}{\sqrt{\bar\alpha_t}}\Bigl(x_t + (1-\bar\alpha_t)\,\nabla_{x_t}\log p_t(x_t)\Bigr).
\end{equation}
Thus, predicting $x_0$ is equivalent (up to scaling) to learning the score.
Score-based diffusion models directly parametrize $s_\theta(x_t,t)\approx \nabla_{x_t}\log p_t(x_t)$.

\begin{remark}[Information-theoretic viewpoint]
Each forward step $x_{t-1}\to x_t$ is a noisy channel that decreases $I(X_0;X_t)$.
The reverse model learns to invert these channels approximately.
This perspective is useful when connecting diffusion models to rate--distortion and successive refinement.
\end{remark}
